% 原书物理页=印刷页31--32；1.9 全节.（定理 1.9.2 之 ii) 的证明与结尾的注记实际印于原书第 33 页顶部、1.10 节标题之前，属本节内容，一并译出附于末尾.）
\section{Lie 群与 Lie 代数之间的对应}\label{sec:lie-1-9}

这里要研究的是从 Lie 群范畴到 Lie 代数范畴的\newterm{lie-functor}{函子}{functor} $G\mapsto\operatorname{Lie}(G)$ 的性质. 我们可以提出的自然问题是：
\begin{enumerate}[label=\roman*)]
\item 每个 $\mathbb R$-Lie 代数都是某个 Lie 群的 Lie 代数吗？
\item $\operatorname{Lie}(G)$ 的每个子 $\mathbb R$-Lie 代数都是 $G$ 的某个 Lie 子群的 Lie 代数吗？
\item 对映射 $\operatorname{Hom}_{Gp.Lie}(G,G')\longrightarrow\operatorname{Hom}_{\mathbb R-AL}(\operatorname{Lie}(G),\operatorname{Lie}(G'))$ 能说些什么？
\end{enumerate}
对 i) 的回答是肯定的，但我们只部分地证明它. 对 ii) 的回答是否定的，正如下面的例子所表明的：取 $G=S^{1}\times S^{1}$，而 $\mathfrak h\subset\operatorname{Lie}(G)=\mathbb R^{2}$ 为一条斜率为无理数的直线. 这条直线可以积分成一个浸入的 Lie 群（直线在环面上的缠绕），却不能积分成一个 Lie 子群. 这个例子提示了 ii) 的一个减弱的版本，而它被证明是正确的：

\begin{theorem}\label{thm:lie-1-9-1}
设 $G$ 是一个 Lie 群，$\mathfrak h\subset\operatorname{Lie}(G)$ 是一个 Lie 子代数. 那么存在一个浸入 Lie 群 $i:H\hookrightarrow G$，其 Lie 代数为 $\mathfrak h$.
\end{theorem}

\begin{remark*}
有了这个定理，问题 i) 就归结为判断每个 Lie 代数是否都能嵌入某个 $\mathfrak{gl}_{n}(\mathbb R)$ 之中. 答案是肯定的（Ado 定理），但我们不打算证明它. 尽管如此我们还是注意到，如果 $\mathfrak g$ 的中心 $\mathfrak z:=\{X\in\mathfrak g,\ [X,\mathfrak g]=0\}$ 是平凡的，那么伴随表示 $\mathfrak g\longrightarrow\mathfrak{gl}(\mathfrak g)$ 就给出了这样一个嵌入.
\end{remark*}

\begin{proof}
这是微分几何中的 Frobenius 定理的一个应用.
\emph{回顾.}\ 设 $M$ 是一个维数为 $m$ 的光滑流形，并设 $n\leqslant m$. 假设给定了由 $\mathbb R$-向量子空间 $D(x)\subset T_{x}M$ 组成、维数为 $n$ 的一族 $D=(D(x))_{x\in M}$. 我们要问的是，是否存在 $M$ 的子流形 $N$ 满足性质
\[
(*)\quad \forall y\in N,\ T_{y}N=D(y).
\]
Frobenius 定理回答了这个问题，而主要工具是
\[
\mathcal V(D)=\{\xi\in\mathcal V(M),\ \forall x\in M,\ \xi(x)\in D(x)\}.
\]
我们首先需要 $D$ 的一个\emph{正则性}性质：如果可以用一些开集 $U$ 覆盖 $M$，使得对其中每个 $U$ 都存在 $\xi^{U}_{1},\cdots,\xi^{U}_{n}\in\mathcal V(D)$，使得对所有 $x\in U$，$D(x)$ 由 $\xi^{U}_{1}(x),\cdots,\xi^{U}_{n}(x)$ 生成，那么就称 $D$ 是一个 $n$-\newterm{lie-distribution}{分布}{distribution}.
\begin{theorem*}
\textup{(Frobenius)}\ 设 $D$ 是 $M$ 上的一个 $n$-分布. 假设 $\mathcal V(D)$ 在向量场的括号下稳定. 那么：
\begin{itemize}[label=--]
\item （存在性）过每一点 $x\in M$ 都有一个满足 $(*)$ 的子流形 $N$.
\item （局部唯一性）若 $N'$ 满足同样的性质，则存在 $x$ 的一个开邻域 $U$，使得 $U\cap N=U\cap N'$.
\end{itemize}
\end{theorem*}
由局部唯一性性质，过点 $x$ 且满足 $(*)$ 的 $M$ 的所有\emph{连通}子流形之并是一个连通的浸入子流形 $N_{x,D}\hookrightarrow M$. $N_{x,D}$ 上的流形结构按如下方式给出. 首先注意到，对如定理中的 $N,N'$，$N\cap N'$ 是 $N$ 与 $N'$ 的开子流形. 然后在 $N_{x,D}$ 上赋予\newterm{lie-inductive-limit-topology}{归纳极限拓扑}{inductive limit topology}. 那么，每个 $N$ 都是 $N_{x,D}$ 的开集，因而 $N_{x,D}$ 上存在唯一的流形结构，使得每个 $N$ 都是它的一个开子流形.
\emph{应用.}\ 设 $\mathfrak h\subset\operatorname{Lie}(G)$ 如定理~\ref{thm:lie-1-9-1} 中所述. 考虑由 $D_{\mathfrak h}(x):=d_{e}L_{x}(\mathfrak h)$ 定义的一族 $D_{\mathfrak h}$. 若 $X_{1},\cdots,X_{n}$ 是 $\mathfrak h$ 的一组基，则对一切 $x\in G$，$D_{\mathfrak h}(x)$ 由 $\xi_{X_{1}}(x),\cdots,\xi_{X_{n}}(x)$ 生成，因此 $D_{\mathfrak h}$ 是 $G$ 上的一个分布，且
\[
\mathcal V(D_{\mathfrak h})=\mathcal C^{\infty}(G)\xi_{X_{1}}\oplus\cdots\oplus\mathcal C^{\infty}(G)\xi_{X_{n}}.
\]
对 $1\leqslant i,j\leqslant n$，有 $[\xi_{X_{i}},\xi_{X_{j}}]=\xi_{[X_{i},X_{j}]}\in\mathbb R\xi_{X_{1}}\oplus\cdots\oplus\mathbb R\xi_{X_{n}}\subset\mathcal V(D_{\mathfrak h})$. 但此时，公式
\[
[f\xi,f'\xi']=ff'[\xi,\xi']+f\partial_{\xi}(f').\,\xi'-f'\partial_{\xi'}(f).\,\xi,\quad\forall f,f'\in\mathcal C^{\infty}(G)
\]
表明 $\mathcal V(D_{\mathfrak h})$ 在向量场的括号下稳定.
设 $H_{g,D_{\mathfrak h}}\hookrightarrow G$ 为由 Frobenius 定理给出的、过点 $g$ 的极大连通浸入子流形. 简单地记 $H:=H_{e,D_{\mathfrak h}}$. 由于 $D_{\mathfrak h}$ 在左平移下不变，故对一切 $g\in G$ 有 $H_{g,D_{\mathfrak h}}=L_{g}(H)$. 而对 $h\in H$，由唯一性还有 $H_{h,D_{\mathfrak h}}=H$. 由此推出 $H$ 对乘法稳定. 此外，$L_{h^{-1}}(H)=L_{h^{-1}}(H_{h,D_{\mathfrak h}})=H$，故 $h^{-1}=L_{h^{-1}}(e)\in H$，即 $H$ 对求逆也稳定. 由于 $H$ 的乘法与求逆的光滑性由 $G$ 的相应光滑性得出，我们得出结论：$H$ 是浸入于 $G$ 的一个 Lie 群. 由构造有 $\operatorname{Lie}(H)=D_{\mathfrak h}(e)=\mathfrak h$.
\end{proof}

现在转向上面提出的问题 iii).

\begin{theorem}\label{thm:lie-1-9-2}
设 $G,G'$ 为两个 Lie 群. 映射
\[
\begin{array}{rcl}
\operatorname{Hom}_{Gp.Lie}(G,G') & \longrightarrow & \operatorname{Hom}_{\mathbb R-AL}\bigl(\operatorname{Lie}(G),\operatorname{Lie}(G')\bigr)\\[2pt]
\varphi & \longmapsto & d\varphi
\end{array}
\]
满足：
\begin{enumerate}[label=\roman*)]
\item 若 $G$ 连通，则它是单射；
\item 若 $G$ 单连通，则它是双射.
\end{enumerate}
\end{theorem}

\begin{proof}
i) 我们知道，对一切 $X\in\operatorname{Lie}(G)$ 与任何态射 $\varphi:G\longrightarrow G'$，有 $\varphi(\exp_{G}(X))=\exp_{G'}(d\varphi(X))$. 于是，若 $d\varphi=d\varphi'$，则 $\varphi$ 与 $\varphi'$ 在 $\exp_{G}$ 的像上重合. 而该像生成连通的 $G$，故 $\varphi=\varphi'$.
ii) 设 $f:\operatorname{Lie}(G)\longrightarrow\operatorname{Lie}(G')$ 是一个 Lie 代数态射. 考虑图像 $\Gamma_{f}:\operatorname{Lie}(G)\longrightarrow\operatorname{Lie}(G)\times\operatorname{Lie}(G')$ 的像 $\mathfrak h$. 根据定理~\ref{thm:lie-1-9-1}，存在一个连通 Lie 群 $H$ 与一个浸入 $\iota:H\hookrightarrow G\times G'$，它诱导一个同构 $d\iota:\operatorname{Lie}(H)\xrightarrow{\ \sim\ }\mathfrak h$. 映射 $\pi:=pr_{1}\circ\iota:H\longrightarrow G$ 是一个 Lie 群态射，其导数 $d\pi=pr_{1}\circ d\iota$ 由构造是同构. 由于 $G$ 连通，$\pi$ 是一个覆盖. 但由于 $G$ 单连通，$\pi$ 是平凡的，因而是一个同构. 此时令 $\varphi:=pr_{2}\circ\iota\circ\pi^{-1}$. 我们有 $d\varphi=f$.
\end{proof}

\begin{remark*}
把存在性定理（它回答了问题 i)）、万有覆盖的存在性定理与上面的定理合在一起，我们得到：\emph{函子 $G\mapsto\operatorname{Lie}(G)$ 是一个\newterm{lie-equivalence-of-categories}{范畴等价}{equivalence of categories}}
\[
\{\text{单连通 Lie 群}\}\xrightarrow{\ \sim\ }\{\text{有限维 $\mathbb R$-Lie 代数}\}.
\]
此外，\emph{该函子把短正合列映为短正合列，且任何\newterm{lie-quasi-inverse}{拟逆}{quasi-inverse}都具有同样的性质}. 这里，一个序列
\[
1\longrightarrow H\xrightarrow{\ \varphi\ }G\xrightarrow{\ \varphi'\ }K\longrightarrow 1
\]
称为``短正合''的，如果 $\varphi$ 是一个嵌入，满足 $\varphi(H)=\operatorname{Ker}(\varphi')$，且 $\varphi'$ 是满射（此时 $\varphi'$ 诱导一个同构 $G/\varphi(H)\xrightarrow{\ \sim\ }K$）. 另一方面，一个序列
\[
(*)\quad 0\longrightarrow\mathfrak h\xrightarrow{\ f\ }\mathfrak g\xrightarrow{\ f'\ }\mathfrak k\longrightarrow 0
\]
称为短正合的，如果它作为向量空间序列是短正合的（此时它诱导一个同构 $\mathfrak g/f(\mathfrak h)\xrightarrow{\ \sim\ }\mathfrak k$）. 由迄今为止已经看到的内容，显然函子 $\operatorname{Lie}$ 把群的短正合列映为 Lie 代数的短正合列. 反过来，从一个 Lie 代数的正合列 $(*)$ 出发，把它（以本质唯一的方式）积分为一个由\emph{单连通} Lie 群组成的序列 $H\xrightarrow{\ \varphi\ }G\xrightarrow{\ \varphi'\ }K$. 那么可以证明，这个群的序列是短正合的. 不显然之点在于 $\varphi$ 是一个嵌入这一事实，它需要用到 Lie 群的 $\pi_{2}$ 是平凡的这一事实.
\end{remark*}
